\documentclass[svgnames]{article}     % use "amsart" instead of "article" for AMSLaTeX format
%\geometry{landscape}                 % Activate for rotated page geometry

%\usepackage[parfill]{parskip}        % Activate to begin paragraphs with an empty line rather than an indent

\usepackage{graphicx}                 % Use pdf, png, jpg, or eps§ with pdflatex; use eps in DVI mode

%maths                                % TeX will automatically convert eps --> pdf in pdflatex
\usepackage{amssymb}
\usepackage{amsmath}
\usepackage{esint}
\usepackage{geometry}

% Inverting Color of PDF
%\usepackage{xcolor}
%\pagecolor[rgb]{0.19,0.19,0.19}
%\color[rgb]{0.77,0.77,0.77}

%noindent
\setlength\parindent{0pt}

%pgfplots
\usepackage{pgfplots}

%images
%\graphicspath{{ }}                   % Activate to set a image directory

%tikz
\usepackage{pgfplots}
\pgfplotsset{compat=1.15}
\usepackage{comment}
\usetikzlibrary{arrows}
\usepackage[most]{tcolorbox}

%Figures
\usepackage{float}
\usepackage{caption}
\usepackage{lipsum}


\title{Stellar Fundamental Metallicity Relation}
\author{Deval Deliwala}
%\date{}                              % Activate to display a given date or no date

\begin{document}
\maketitle
%\section{}
%\subsection{}
%\tableofcontents                     % Activate to display a table of contents

\section{Background Information}


The gas in a galaxy's interstellar medium (ISM) and the stars themselves evolve
together over time, with their chemical compositions influencing each other. \\

The metallicity (fraction of heavier elements relative to all baryonic matter)
of a galaxy reflects a history of star formation, gas inflow and outflow,
throughout time. As stars form and explode, they enrich the ISM with metals.
Along with other galactic events like mergers or gas inflow/outflow,
influences the amount and distribution of metals, encoding a history of the
galaxies' interactions and internal processes.\\

Previous studies show that stellar and gas metallicities scale with stellar
mass, known as the \textit{mass-metallicity} relations (MZRs). The more massive
galaxies on average are more chemically enriched than lower mass galaxies. Gas
metallicity has also been shown to depend on the recent star formation
activity of a galaxy (gFMR). When a galaxy experiences bursts of star formation, they
use up gas and may lead to higher metallicity as supernovae eject metals into
the surrounding gas. However, intense star formation can also reduce
metallicity if the enriched gas is expelled entirely from the galaxy. 

\subsection{Previous Results}

\begin{itemize}
  \item At a fixed stellar mass, star-forming galaxies tend to have
    \textit{lower} metallicities than quiescent, inactive galaxies. 
  \item Galaxies in between star-forming and inactive phases follow an
    mass metallicity  relation in the middle 
  \item Galaxies with younger stars, like spiral galaxies, generally have lower
    metallicities than older galaxies at \textit{all} mass levels. 
\end{itemize}

\paragraph{The Goal of the Study} \mbox{} \\

This study extends previous research by exploring the stellar mass-metallicity
relation (sFMR) in a new way -- focusing on recent star formation activity.
Instead of solely looking at stellar ages, they analyze how stellar metallicity
varies with a galaxy's position relative to the \textit{Star-Forming Main
Sequence} (SFMS) -- which is statistical relationship showing how galaxy's
star formation rates (SFR) scale with their stellar masses. They rename this
relation as the sFMR -- or \textit{stellar fundamental metallicity relation},
which links metallicity  to \textit{both} stellar mass and current star
formation. This is to show that stellar metallicity in galaxies can also be
understood via the combined effects of stellar mass and recent star formation
activity. 

\section{Collecting Data}

They use the MaNGA dataset from the SDSS DR17. This data has spatially resolved
spectra, a very high signal-to-noise ratio and a large sample size. To compare
their MaNGA results with current cosmological simulations, they collect an
analogous sample of simulated galaxies from the IlustricTNG simulation suite
that contains stellar masses, mass-weighted mean metal mass fractions of star
particles.


\section{Procedure}

They begin by calculating distance from the SFMS for every galaxy in the MaNGA
sample. This distance quantifies how a galaxy's star formation rate compares to
the expected value for a galaxy of its mass. 

\paragraph{Figure 1 \& 2} \mbox{} \\ 

The stellar mass - star formation rate diagram, which visually represents the
relationship between galaxy's stellar mass and its star formation rate (SFR).
\\

A prominent feature of this diagram is the SFMS -- the linear relationship
between the total stellar mass and the SFR for star-forming galaxies. Here it
is represented by the white dashed-line. Actively star-forming galaxies lie
along this sequence. They calculate a best-fit of the SFMS as 

\[
  MS(M_*) = -7.96 + 0.76\log_{10}(M_*[M_\odot]).
\] \vspace{3px}

When doing the same calculation with their synthetic galaxy sample from
IllustrisTNG, they calculated a value of 

\[
  MS_\text{TNG}(M_*) = -6.28 + 0.63 \log_{10} (M_* [M_\odot]). 
\] \vspace{3px}

\begin{itemize}
  \item Any spatial pixel with an $SNR>1$ is discarded. Afterwards, the data is
    binned in radial annuli to ensure all the data in each bin has a higher
    median SNR. 
  \item The spectra from each spaxel is summed, resulting in a binned spectra
    for analysis. 
  \item The binned spectra is then fitted using a customized $\chi^2$ 
    minimization Penalized Pixel-Fitting (PPXF) code which accounts for both
    the stellar and gas components of the spectra 
    \begin{itemize}
      \item \textbf{Gas Emission Lines}: Fitted using Gaussians that describe
        the distribution of light from the emission lines 
      \item \textbf{Stellar Continuum}: Fitted using a library of Simple
        Stellar Population (SSP) templates which account for various ages and
        metallicities
    \end{itemize}
  \item This method is referred to as  \textit{astro-archaeological} approach
    because it treats the observed stellar populations as \textit{fossil
    records} of a galaxy's star formation history. 
  \item By fitting the spectra with both SSPs and gas emission lines, they say
    they can disentangle the contributions of different stellar populations and
    their star formation activity -- helping reconstruct the galaxy's history,
    including periods of intense star formation and inactivity. 
\end{itemize}

\paragraph{More Detailed Steps} \mbox{} \\

\begin{itemize}
  \item Step 1 -- Sky Masking \& Normalization
    \begin{itemize}
      \item Mask the sky emission lines to reduce atmospheric interference.
      \item Normalize the observed spectrum and SSP templates by the median
        flux per  spectral pixel, preventing overfitting to noise in teh data. 
    \end{itemize}
  \item Step 2 -- Initial PPXF Fit
    \begin{itemize}
      \item An initial PPXF fit is performed to estimate the intrinsic noise in
        the spectrum and filter out problematic spectral pixels 
    \end{itemize}
  \item Step 3 -- $3\sigma$ Clipping
    \begin{itemize}
      \item a $3\sigma$ clipping processes is performed. Gas emission lines are
        excluded since they are allowed to have really high fluxes and thus
        high residuals 
    \end{itemize}
  \item Step 4 -- Refitting with PPXF
    \begin{itemize}
      \item The spectrum is fit again with PPXF using the updated noise
        estimate, yielding a new best-fit solution $y(\lambda)$. 
    \end{itemize}
  \item Step 5 -- Residual-Based Bootstrapping
    \begin{itemize}
      \item A bootstrapping method is employed to assess the fits reliability.
        The best-fit solution $y(\lambda)$ is perturbed by the residuals from
        the initial fit, leading to a modified spectrum $y^*(\lambda)$. This
        perturbation involves adding and removing random residuals from the
        original fit. 
    \end{itemize}
  \item Step 6 -- Iterative Fitting
    \begin{itemize}
      \item The perturbed spectrum is again fit using PPXF, without any noise
        estimate. 
    \end{itemize}
  \item Averaging over Iterations
    \begin{itemize}
      \item Steps 5-6 are repeated 100 times averaging all the age-metallicity
        weights to recover the star-formation history consistent with the
        intrinsic noise of the spectrum. 
    \end{itemize}
\end{itemize}

This process does not depend on any assumption about the underlying physics of
galaxy evolution. 


\section{The Stellar Fundamental Metallicity Relation}

\paragraph{Figure 3} \mbox{} \\ 

Here they present the global sFMR. By summing over the spectral analysis
results over spatially resolved measurements, they find a fit for 6 different
$\Delta_{MS}$ bins. \\

The positive slope indicates that as the stellar mass of a galaxy increases,
its light-weighted average metallicity increases. The more massive galaxies
tend to have higher metallicities. \\

Additionally, as $\Delta_{MS}$ increases (galaxies are more star-forming),
their metallicity decreases relative to the inactive ones. 

This also suggests more massive galaxies have more efficient processes for
converting gas into stars and recycling them back into the ISM. \\

\paragraph{Figure 5} \mbox{} \\ 

They repeat the analysis but focusing on younger stellar populations
$<10^{8.5}$ yrs. The found that more star-forming galaxies have \textit{lower}
metallicities than the inactive populations across all stellar populations. The
young metallicity values however are systematically higher than the total
metallicity values for galaxies at all masses and distances from the MS.


\section{Starvation Hypothesis}

The difference in metallicity between star-forming (SF) and quiescent galaxies
is interpreted as evidence for a phenomenon called starvation. This hypothesis
posits that even after the supply of low-enrichment gas ceases, star formation
can continue in a galaxy. \\

As a result, the ISM (interstellar medium) becomes increasingly enriched with
metals because the same gas is recycled through multiple star formation
episodes without being diluted by newly accreted low-metallicity gas. \\

The underlying assumption is that galaxies with low-metallicity/high star
formation rates (MS) will evolve into high-metallicity/low star formation rate
(quiescent) galaxies over time.\\ 

Comparing different MS bins in the local Universe may not capture this
evolutionary trend, as these galaxies are not necessarily linked by direct
progenitor–descendant relationships. However, if the metallicities of quiescent
galaxies were compared to their SF progenitors, the metallicity gap would be
even more pronounced. 

\section{Metal Retention -- Gravitational Potential Wells} 

In contrast, the Gravitational Well hypothesis posits that the difference in
MZR between SF and inactive galaxies is because galaxies with a higher
$M_*/R_e$ c an retain more metals as their gravitational potential reduces the
outflows of metal enriched gas. \\

However, they implemented this idea to estimate how much outflow would occur,
but found a small metallicity offset to still persist, implying that even
though it does occur, other processes are likely more dominant in forming the
sFMR. 

\section{Conclusion}

The gas-phase Fundamental Metallicity Relation (gFMR) connects star formation
and gas metallicity over long time scales, while the stellar FMR (sFMR) extends
this correlation, linking the star formation rate (SFR) to the light-weighted
metallicity of the stars over similar time frames.\\

This relationship suggests that galaxies experience continuous inflows of
low-metallicity gas from the intergalactic medium (IGM) or circumgalactic
medium (CGM), which not only fuels star formation but also maintains a balance
between metal production and dilution in the ISM.\\

The correlation observed for SF galaxies suggests that the processes governing
their metallicity and star formation operate over extended periods rather than
being driven by short bursts of star formation or gas accretion. \\

This aligns with theoretical models predicting long-lasting gas inflow
mechanisms that regulate star formation and chemical enrichment in galaxies.

\end{document}

